\subsection{Impact de l'EMA sur l'apprentissage}\label{sub:Impact de l'EMA sur l'apprentissage} % (fold)

La Figure~\ref{fig:ema_dynamics} illustre comment le momentum $m$
contrôle la vitesse de convergence du professeur vers l'élève.
 
% ------------------------------------------------------------------
% Figure 4 : dynamique EMA selon m
% ------------------------------------------------------------------
\begin{figure}[ht]
\centering
\begin{tikzpicture}
\begin{axis}[
  width=0.72\linewidth, height=4.8cm,
  xlabel={Itération $t$},
  ylabel={$\|\xi_t - \theta_t\|_2$ (normalisé)},
  xmin=0, xmax=300,
  ymin=0, ymax=1.05,
  xtick={0,50,100,150,200,250,300},
  ytick={0,0.25,0.5,0.75,1.0},
  grid=both, grid style={line width=0.3pt, draw=gray!25},
  axis background/.style={fill=gray!4},
  legend style={at={(0.97,0.97)}, anchor=north east,
                font=\tiny, fill=white, draw=gray!40},
  label style={font=\small},
  tick label style={font=\tiny},
]
% m = 0.90
\addplot[colMask, thick, smooth] coordinates {
  (0,1)(10,0.65)(20,0.42)(40,0.18)(60,0.08)(100,0.02)(200,0.001)(300,0)
};
\addlegendentry{$m=0{,}90$}
% m = 0.996 (notre valeur)
\addplot[colCtx, thick, smooth, dashed] coordinates {
  (0,1)(20,0.92)(50,0.82)(100,0.67)(150,0.55)(200,0.44)(300,0.29)
};
\addlegendentry{$m=0{,}996$ (utilisé)}
% m = 0.9999
\addplot[colTgt, thick, smooth, dotted] coordinates {
  (0,1)(50,0.995)(100,0.99)(200,0.98)(300,0.97)
};
\addlegendentry{$m=0{,}9999$}
\end{axis}
\end{tikzpicture}
\caption{Évolution schématique de la distance $\|\xi_t -
  \theta_t\|_2$ pour différentes valeurs de momentum $m$.
  Un $m$ trop faible (0{,}90) rend le professeur trop réactif et
  instabilise les cibles. Un $m$ trop élevé (0{,}9999) ralentit
  excessivement la convergence. La valeur $m=0{,}996$ offre le
  meilleur compromis pour notre échelle d'entraînement.}
\label{fig:ema_dynamics}
\end{figure}
 
% subsection Impact de l'EMA sur l'apprentissage (end)
